\documentclass[a4paper,12pt]{article}

\usepackage[T2A]{fontenc}
\usepackage[utf8]{inputenc}
\usepackage[russian]{babel}
\usepackage{amsmath,amssymb,mathtools}
\usepackage[scaled=0.95]{DejaVuSans}
\renewcommand{\familydefault}{\sfdefault}
\usepackage{icomma}
\usepackage[a4paper,margin=20mm,headheight=25pt]{geometry}
\usepackage[nopatch=footnote]{microtype}
\usepackage{tikz}
\usetikzlibrary{calc,arrows.meta,positioning,shapes.geometric}
\usepackage{pgfplots}
\pgfplotsset{compat=1.18}
\usepackage{tabularx,booktabs}
\usepackage{enumitem}
\usepackage{hyperref}
\usepackage{parskip}
\usepackage{fancyhdr}
\usepackage{setspace}
\usepackage{xcolor}

% ------------------------------------------------
% Палитра и базовая типографика
% ------------------------------------------------
\definecolor{accent}{HTML}{008080}
\definecolor{darkaccent}{HTML}{004D4D}
\definecolor{textgray}{HTML}{333333}
\definecolor{softgray}{HTML}{666666}

\color{textgray}
\onehalfspacing
\setlength{\parindent}{0pt}
\setlength{\parskip}{0.5em}
\setlength{\emergencystretch}{2em}

\hypersetup{
    hidelinks,
    pdftitle={Чек-лист: рациональные функции и гипербола},
    pdfauthor={Лёвин Артём Александрович}
}

% ------------------------------------------------
% Колонтитулы
% ------------------------------------------------
\pagestyle{fancy}
\fancyhf{}
\fancyhead[L]{\small Марина Селиверстова}
\fancyhead[R]{\small ОГЭ --- математика}
\fancyfoot[C]{\small\thepage}
\renewcommand{\headrulewidth}{0.4pt}
\renewcommand{\headrule}{%
  \hbox to\headwidth{\color{accent}\leaders\hrule height \headrulewidth\hfill}}

% ------------------------------------------------
% Списки и заголовки
% ------------------------------------------------
\setlist[itemize]{
    leftmargin=1.7em,
    itemsep=0.25em,
    topsep=0.25em
}

\setlist[enumerate]{
    leftmargin=2em,
    itemsep=0.35em,
    topsep=0.25em
}

\newcommand{\sectiontitle}[1]{%
    \par\vspace{0.65em}
    {\Large\bfseries\color{darkaccent}#1}\par
    \vspace{0.15em}
    {\color{accent}\rule{\linewidth}{0.65pt}}\par
    \vspace{0.3em}
}

\newcommand{\subsectiontitle}[1]{%
    \par\vspace{0.45em}
    {\large\bfseries\color{darkaccent}#1}\par
    \vspace{0.1em}
}

\newcommand{\checkitem}{\(\square\)\;}

% ------------------------------------------------
% TikZ-стили блок-схемы
% ------------------------------------------------
\tikzset{
    flowstart/.style={
        ellipse,
        draw=darkaccent,
        line width=0.9pt,
        align=center,
        minimum width=38mm,
        minimum height=10mm,
        font=\small
    },
    flowaction/.style={
        rectangle,
        rounded corners=3pt,
        draw=accent,
        line width=0.9pt,
        align=center,
        text width=120mm,
        minimum height=12mm,
        inner sep=4.5pt,
        font=\small
    },
    flowdecision/.style={
        diamond,
        aspect=2.35,
        draw=darkaccent,
        line width=0.9pt,
        align=center,
        text width=58mm,
        inner sep=2.5pt,
        font=\small
    },
    flowbranch/.style={
        rectangle,
        rounded corners=3pt,
        draw=accent,
        line width=0.9pt,
        align=center,
        text width=62mm,
        minimum height=12mm,
        inner sep=4pt,
        font=\small
    },
    flowarrow/.style={
        -{Stealth[length=3mm,width=2mm]},
        line width=0.9pt,
        draw=textgray
    }
}

% ------------------------------------------------
% Локальная A3-страница для крупной блок-схемы
% ------------------------------------------------
\newenvironment{flowchartpageportrait}{%
    \clearpage
    \begingroup
    \setlength{\paperwidth}{297mm}%
    \setlength{\paperheight}{420mm}%
    \pdfpagewidth=297mm
    \pdfpageheight=420mm
    \setlength{\hoffset}{0pt}%
    \setlength{\voffset}{0pt}%
    \setlength{\oddsidemargin}{-1.4mm}%
    \setlength{\evensidemargin}{-1.4mm}%
    \setlength{\topmargin}{-3.4mm}%
    \setlength{\headheight}{0pt}%
    \setlength{\headsep}{0pt}%
    \setlength{\footskip}{0pt}%
    \setlength{\textwidth}{249mm}%
    \setlength{\textheight}{376mm}%
    \setlength{\linewidth}{249mm}%
    \setlength{\columnwidth}{249mm}%
    \hsize=249mm
    \vsize=376mm
    \thispagestyle{empty}%
}{%
    \clearpage
    \endgroup
}

\begin{document}

% ================================================================
% ТИТУЛЬНЫЙ ЛИСТ
% ================================================================
\thispagestyle{empty}

\begin{center}

\vspace*{1.4cm}

{\large\color{softgray} Подготовка к ОГЭ по математике}

\vspace{1.4cm}

{\fontsize{28}{34}\selectfont\bfseries\color{darkaccent}
Чек-лист}

\vspace{0.55cm}

{\LARGE\bfseries
Рациональные функции и гипербола}

\vspace{0.55cm}

{\large
асимптоты \(\boldsymbol{\cdot}\)
точки пересечения \(\boldsymbol{\cdot}\)
ОДЗ\\[0.15em]
сокращение дробей \(\boldsymbol{\cdot}\)
выколотые точки \(\boldsymbol{\cdot}\)
прямая \(y=m\)}

\vspace{1.45cm}

{\Large\bfseries Марина Селиверстова}

\vspace{0.35cm}

{\large Тренировочный блок: Путь А}

\vfill

{\large \today}

\vspace{1.45cm}

{\small
Лёвин Артём Александрович\\
эксклюзивно для Марины Селиверстовой}

\vspace{0.9cm}

\begin{tikzpicture}[remember picture,overlay]
    \draw[
        accent,
        line width=1.15pt
    ]
    ([xshift=18mm,yshift=18mm]current page.south west)
    .. controls
    ([xshift=63mm,yshift=29mm]current page.south west)
    and
    ([xshift=-70mm,yshift=11mm]current page.south east)
    ..
    ([xshift=-18mm,yshift=20mm]current page.south east);
\end{tikzpicture}

\end{center}

\clearpage

% ================================================================
% КЛЮЧЕВЫЕ ЗНАНИЯ
% ================================================================
\sectiontitle{1. Что необходимо уметь}

После повторения темы Вы должны уверенно выполнять следующие действия.

\begin{itemize}
    \item \checkitem определять ОДЗ исходной рациональной функции;
    \item \checkitem упрощать рациональное выражение, не теряя исходные ограничения;
    \item \checkitem находить вертикальную и горизонтальную асимптоты;
    \item \checkitem находить точки пересечения графика с осями \(Ox\) и \(Oy\);
    \item \checkitem определять расположение ветвей гиперболы;
    \item \checkitem отмечать выколотые точки после сокращения множителей;
    \item \checkitem использовать прямую \(y=m\) для исследования числа общих точек с графиком.
\end{itemize}

\sectiontitle{2. Базовая модель гиперболы}

Удобнейший вид для построения:

\[
\boxed{y=a+\frac{b}{x-c}},
\qquad b\ne0.
\]

Из формулы сразу читаются асимптоты:

\[
\boxed{x=c},
\qquad
\boxed{y=a}.
\]

Точка

\[
\boxed{C(c;a)}
\]

--- центр симметрии графика, то есть точка пересечения асимптот.

\subsectiontitle{Как знак \(b\) влияет на ветви}

Если

\[
b>0,
\]

ветви расположены относительно центра \(C\) так же, как у графика

\[
y=\frac1x.
\]

Если

\[
b<0,
\]

ветви расположены относительно центра \(C\) так же, как у графика

\[
y=-\frac1x.
\]

\textbf{Важно:} в общем случае симметрия выполняется относительно точки
\(C(c;a)\), а не относительно начала координат.

\sectiontitle{3. Линейно-дробная функция}

Пусть

\[
y=\frac{Ax+B}{Cx+D},
\qquad C\ne0.
\]

Если

\[
AD-BC\ne0,
\]

график является гиперболой с асимптотами

\[
\boxed{x=-\frac{D}{C}},
\qquad
\boxed{y=\frac{A}{C}}.
\]

Если \(AD-BC=0\), числитель и знаменатель пропорциональны. После упрощения
может получиться постоянная функция с исключённой точкой, поэтому формулу
гиперболы применять нельзя механически.

\subsectiontitle{Пример}

Для функции

\[
y=\frac{x+1}{x-1}
\]

получаем

\[
x=1,
\qquad
y=1
\]

как вертикальную и горизонтальную асимптоты соответственно.

Точки пересечения с осями:

\[
y=0
\Longrightarrow
x+1=0
\Longrightarrow
\boxed{(-1;0)},
\]

\[
x=0
\Longrightarrow
y=-1
\Longrightarrow
\boxed{(0;-1)}.
\]

% ================================================================
% ОДЗ И СОКРАЩЕНИЕ
% ================================================================
\sectiontitle{4. ОДЗ и сокращение рациональной дроби}

ОДЗ определяют \textbf{по исходной формуле до сокращения}.

Рассмотрим

\[
y=\frac{x+1}{x^2+x}-2.
\]

Разложим знаменатель:

\[
x^2+x=x(x+1).
\]

Следовательно,

\[
\boxed{x\ne0,\quad x\ne-1}.
\]

Теперь можно сократить общий множитель:

\[
y=
\frac{x+1}{x(x+1)}-2
=
\frac1x-2,
\qquad x\ne0,-1.
\]

Ограничение \(x\ne-1\) после сокращения сохраняется. При \(x=-1\)
упрощённая формула дала бы

\[
y=-1-2=-3.
\]

Поэтому точка

\[
\boxed{(-1;-3)}
\]

на графике должна быть \textbf{выколота}.

\subsectiontitle{Правило}

Если

\[
\frac{(x-a)P(x)}{(x-a)Q(x)}
\]

сокращается на \(x-a\), то после сокращения обязательно остаётся условие

\[
\boxed{x\ne a}.
\]

Сокращение изменяет запись функции, но не возвращает в её область определения
значения, запрещённые в исходной формуле.

\sectiontitle{5. Точки пересечения с координатными осями}

\subsectiontitle{Пересечение с \(Ox\)}

На оси \(Ox\)

\[
y=0.
\]

Для дроби

\[
\frac{P(x)}{Q(x)}=0
\]

нужно одновременно выполнить условия

\[
\boxed{P(x)=0},
\qquad
\boxed{Q(x)\ne0}.
\]

\subsectiontitle{Пересечение с \(Oy\)}

На оси \(Oy\)

\[
x=0.
\]

Поэтому нужно подставить \(x=0\) в функцию. Если \(0\notin D(f)\),
пересечения с \(Oy\) нет.

\sectiontitle{6. Быстрое чтение сдвинутой гиперболы}

Если функция уже записана в виде

\[
y=a+\frac{b}{x-c},
\]

не требуется предварительно приводить её к одной большой дроби.

Например,

\[
y=3+\frac1x.
\]

Сразу получаем

\[
\boxed{x=0},
\qquad
\boxed{y=3}.
\]

Для

\[
y=\frac1{2x+3}-1
\]

горизонтальная асимптота видна сразу:

\[
\boxed{y=-1}.
\]

Вертикальную асимптоту находим из знаменателя:

\[
2x+3=0
\Longrightarrow
\boxed{x=-\frac32}.
\]

Пересечение с \(Ox\):

\[
0=\frac1{2x+3}-1
\Longrightarrow
\frac1{2x+3}=1=\frac11.
\]

После перемножения крест-накрест:

\[
2x+3=1
\Longrightarrow
\boxed{x=-1}.
\]

Пересечение с \(Oy\):

\[
y(0)=\frac13-1=-\frac23,
\]

то есть точка

\[
\boxed{\left(0;-\frac23\right)}.
\]

Кроме того,

\[
y=-1+\frac{\frac12}{x+\frac32},
\]

поэтому коэффициент при дроби положителен и ветви расположены относительно
центра

\[
\left(-\frac32;-1\right)
\]

так же, как у \(y=\frac1x\).

Это быстрый способ. Универсальный способ остаётся тем же: сначала привести
выражение к знакомому виду, затем определить ОДЗ, асимптоты и ключевые точки.

% ================================================================
% ПРИМЕРЫ И ПАРАМЕТР
% ================================================================
\sectiontitle{7. Полный пример: линейно-дробная функция}

Рассмотрим

\[
y=\frac{2x+1}{-x+2}.
\]

ОДЗ:

\[
-x+2\ne0
\Longrightarrow
\boxed{x\ne2}.
\]

Здесь

\[
A=2,
\quad
B=1,
\quad
C=-1,
\quad
D=2.
\]

Асимптоты:

\[
x=-\frac{D}{C}
=-\frac{2}{-1}
=\boxed{2},
\]

\[
y=\frac{A}{C}
=\frac{2}{-1}
=\boxed{-2}.
\]

Пересечение с \(Ox\):

\[
2x+1=0
\Longrightarrow
x=-\frac12,
\]

то есть

\[
\boxed{\left(-\frac12;0\right)}.
\]

Пересечение с \(Oy\):

\[
y(0)=\frac12,
\]

поэтому

\[
\boxed{\left(0;\frac12\right)}.
\]

Этих данных достаточно для качественного эскиза. Большую таблицу значений
строить не требуется.

\sectiontitle{8. Если положение ветвей вызывает сомнение}

Таблица значений остаётся допустимым способом самопроверки.

Удобно брать значения \(x\), расположенные симметрично относительно
вертикальной асимптоты. Например, если асимптота

\[
x=2,
\]

можно взять пары

\[
1\ \text{и}\ 3,
\qquad
0\ \text{и}\ 4.
\]

Обычно достаточно одной-двух дополнительных точек на каждой стороне
асимптоты, чтобы проверить расположение ветвей.

\sectiontitle{9. Выколотая точка и прямая \(y=m\)}

Вернёмся к функции

\[
y=\frac{x+1}{x^2+x}-2.
\]

После упрощения:

\[
y=\frac1x-2,
\qquad
\boxed{x\ne0,-1}.
\]

Асимптоты:

\[
\boxed{x=0},
\qquad
\boxed{y=-2}.
\]

Пересечение с \(Ox\):

\[
0=\frac1x-2
\Longrightarrow
\frac1x=2
\Longrightarrow
x=\frac12,
\]

поэтому

\[
\boxed{\left(\frac12;0\right)}.
\]

Пересечения с \(Oy\) нет, поскольку \(x=0\) не входит в ОДЗ.

При исключённом значении \(x=-1\) упрощённая функция дала бы

\[
y=-3.
\]

Следовательно, точка

\[
\boxed{(-1;-3)}
\]

выколота.

\subsectiontitle{График}

\begin{center}
\begin{tikzpicture}
\begin{axis}[
    width=0.82\linewidth,
    height=0.55\linewidth,
    axis lines=middle,
    xmin=-4,
    xmax=4,
    ymin=-6,
    ymax=4,
    samples=240,
    clip=true,
    unit vector ratio=1 1,
    xlabel={$x$},
    ylabel={$y$},
    xtick={-3,-2,-1,1,2,3},
    ytick={-5,-4,-3,-2,-1,1,2,3},
    grid=both,
    grid style={line width=.2pt,draw=gray!25},
    tick label style={font=\small},
    label style={font=\small}
]

\addplot[
    domain=-4:-1.015,
    thick,
    color=darkaccent
]
{1/x-2};

\addplot[
    domain=-0.985:-0.08,
    thick,
    color=darkaccent
]
{1/x-2};

\addplot[
    domain=0.08:4,
    thick,
    color=darkaccent
]
{1/x-2};

\addplot[
    dashed,
    color=accent,
    thick,
    domain=-4:4
]
{-2};

\addplot[
    only marks,
    mark=o,
    mark size=3.2pt,
    line width=1pt,
    color=darkaccent
]
coordinates {(-1,-3)};

\addplot[
    only marks,
    mark=*,
    mark size=2.2pt,
    color=darkaccent
]
coordinates {(0.5,0)};

\node[anchor=south east,font=\small]
at (axis cs:-1,-3) {$(-1;-3)$};

\node[anchor=south west,font=\small]
at (axis cs:0.5,0) {$\left(\frac12;0\right)$};

\node[anchor=south east,font=\small,text=accent]
at (axis cs:3.8,-2) {$y=-2$};

\node[anchor=south west,font=\small,text=accent]
at (axis cs:0.08,3.1) {$x=0$};

\end{axis}
\end{tikzpicture}
\end{center}

\subsectiontitle{Как работает прямая \(y=m\)}

Прямая

\[
y=m
\]

всегда горизонтальна. Изменение \(m\) перемещает её вверх или вниз.

Для функции

\[
y=-2+\frac1x
\]

значение

\[
\boxed{m=-2}
\]

не даёт общих точек: это уровень горизонтальной асимптоты.

Из-за выколотой точки отсутствует также уровень

\[
\boxed{m=-3}.
\]

Действительно,

\[
\frac1x-2=-3
\Longrightarrow
x=-1,
\]

но \(x=-1\) исключено из ОДЗ. Других решений это уравнение не имеет.

Итак,

\[
\boxed{m\in\{-3;-2\}}.
\]

\sectiontitle{10. Типичные ошибки}

\begin{enumerate}
    \item \textbf{Сократить дробь и забыть исходное ограничение.}

    Исключённое значение остаётся запрещённым после сокращения.

    \item \textbf{Начать строить график до упрощения функции.}

    Сначала определите ОДЗ и приведите выражение к удобному виду.

    \item \textbf{Потерять знак в знаменателе.}

    Например,

    \[
    -x+2=0
    \Longrightarrow
    x=2.
    \]

    \item \textbf{Искать пересечение с \(Oy\), подставляя запрещённое \(x=0\).}

    Если \(0\notin D(f)\), пересечения с \(Oy\) нет.

    \item \textbf{Приравнять дробь к нулю и забыть знаменатель.}

    Для

    \[
    \frac{P(x)}{Q(x)}=0
    \]

    требуется одновременно

    \[
    P(x)=0,
    \qquad
    Q(x)\ne0.
    \]

    \item \textbf{Считать любую гиперболу симметричной относительно начала координат.}

    Центр симметрии графика

    \[
    y=a+\frac{b}{x-c}
    \]

    равен \((c;a)\).

    \item \textbf{Перепутать выколотую точку и асимптоту.}

    Выколотая точка --- одно отсутствующее значение графика. Асимптота описывает
    поведение функции при приближении к определённой прямой.
\end{enumerate}

\clearpage

% ================================================================
% ТРЕНИРОВОЧНЫЙ БЛОК
% ================================================================
\begin{center}
{\LARGE\bfseries\color{darkaccent}
Тренировочный блок}

\vspace{0.2cm}

{\large Путь А}

\vspace{0.15cm}


\end{center}

\vspace{0.4cm}

\sectiontitle{Задание 1}

Постройте график функции

\[
\boxed{y=\frac{2x-3}{x+1}}.
\]

В ходе решения:

\begin{enumerate}[label=\alph*)]
    \item укажите ОДЗ;
    \item найдите вертикальную и горизонтальную асимптоты;
    \item найдите точки пересечения с \(Ox\) и \(Oy\);
    \item постройте эскиз графика.
\end{enumerate}

\vspace{0.6cm}

\sectiontitle{Задание 2}

Постройте график функции

\[
\boxed{y=\frac1{2x+3}-1}.
\]

В ходе решения:

\begin{enumerate}[label=\alph*)]
    \item укажите ОДЗ;
    \item найдите обе асимптоты;
    \item найдите точки пересечения с координатными осями;
    \item определите положение ветвей относительно точки пересечения асимптот;
    \item постройте эскиз графика.
\end{enumerate}

\vspace{0.6cm}

\sectiontitle{Задание 3}

Дана функция

\[
\boxed{y=\frac{x-2}{x^2-2x}-1}.
\]

\begin{enumerate}[label=\alph*)]
    \item Укажите ОДЗ исходной функции.
    \item Упростите выражение, сохранив все ограничения.
    \item Найдите асимптоты.
    \item Найдите точки пересечения с координатными осями.
    \item Найдите и отметьте выколотую точку.
    \item Постройте график.
    \item Найдите все значения \(m\), при которых прямая
    \[
    y=m
    \]
    не имеет с графиком функции ни одной общей точки.
\end{enumerate}

\clearpage

% ================================================================
% ОТВЕТЫ
% ================================================================
\begin{center}
{\LARGE\bfseries\color{darkaccent}
Ответы к тренировочному блоку}
\end{center}

\vspace{0.55cm}

\renewcommand{\arraystretch}{1.4}

\begin{table}[ht]
\centering
\small
\begin{tabularx}{\linewidth}{
    >{\centering\arraybackslash}p{0.06\linewidth}
    >{\raggedright\arraybackslash}X
    >{\raggedright\arraybackslash}X
}
\toprule
\textbf{№} &
\textbf{ОДЗ и асимптоты} &
\textbf{Ключевые точки и итог} \\
\midrule

1 &
\(
D(f):\ x\ne-1.
\)

\(
x=-1,\quad y=2.
\)
&
\(
\left(\frac32;0\right),
\quad
(0;-3).
\)
\\

\midrule

2 &
\(
D(f):\ x\ne-\frac32.
\)

\(
x=-\frac32,\quad y=-1.
\)
&
\(
(-1;0),
\quad
\left(0;-\frac23\right).
\)

Ветви расположены как у \(y=\frac1x\) относительно центра
\(\left(-\frac32;-1\right)\).
\\

\midrule

3 &
\(
D(f):\ x\ne0,\ 2.
\)

\(
y=\frac1x-1,\quad x\ne0,2.
\)

\(
x=0,\quad y=-1.
\)
&
Пересечение с \(Ox\): \((1;0)\).

С \(Oy\) пересечения нет.

Выколотая точка:

\[
\left(2;-\frac12\right).
\]

\[
\boxed{m\in\left\{-1;-\frac12\right\}}.
\]
\\

\bottomrule
\end{tabularx}
\end{table}

\vfill

\begin{center}
{\color{accent}\rule{0.72\linewidth}{0.65pt}}

\vspace{0.3cm}

{\small
Перед построением графика начинайте с ОДЗ и упрощения выражения.\\
Сокращённый множитель исчезает из формулы, но соответствующее ограничение остаётся.}
\end{center}

\end{document}